\subsection{Methods for optimizing \persona{s} in-sample}

We briefly describe two methods for optimizing \persona{s} with a given set of training data.
The first, the selection method, assumes we have a finite library of candidate \persona{s} and selects (or mixes) them to best fit the training data. 
We apply this method to the experiments presented in Section~\ref{sec:predict_new_AR} with data from \cite{1120Arad2012}.
An early version of this idea was suggested by \cite{horton2023large}, and several others have since explored applications \citep{leng2024calibrate,Jackson2025Mixture,bui2025mixture}.
The second, the construction method, parameterizes a prompt template with numeric trait dimensions and optimizes those parameters.
To the best of our knowledge, this method is novel.
An application of the method is presented in Appendix~\ref{app:predict_new_CR} using data from \cite{charness2002understanding}.

\paragraph{Selection Method.}
A finite set of unique candidate natural language \persona{s} is first specified. 
For each \persona{} $\theta \in \Theta$, the LLM is used to generate a predictive distribution $\hat P_\theta$. 
Let $P$ represent the observed ground-truth human distribution.
The objective is to solve
\[
\min_{\mathbf{w}} \; d \Bigl(P, \sum_{\theta \in \Theta} w_\theta \,\hat P_\theta\Bigr) \quad \text{subject to} \quad \sum_{\theta \in \Theta} w_\theta = 1, \quad w_\theta \ge 0,
\]
where $d$ is a chosen distance measure (e.g., KL divergence or the mean absolute distance between distributions). 
Once solved, these weights can be used to scale the appropriate mixture of \persona{s} (i.e., $\bm{\theta^{\star}}$) and applied to new settings.

\paragraph{Construction Method.}
Alternatively, a \persona{} template can be parameterized by numeric trait variables. 
This is best illustrated with an example.
Suppose $\phi_1$ and $\phi_2$ capture degrees of self-interest and inequity aversion, respectively. 
The \persona{} could be:
\begin{quote}
\singlespacing
\vspace{-.5cm}
$\theta(\phi_1, \phi_2)$ = \emph{``You weigh your own payoff with weight $\{\phi_1\}$, and you dislike creating disadvantageous inequality at level $\{\phi_2\}$. Please respond accordingly.''}
\end{quote}
$\hat P_\theta$ denotes the distribution induced by the LLM under parameter vector $\theta = (\phi_1, \phi_2)$. 
Given an observed human distribution $P$, the optimal parameters $\bm{\theta^{\star}}$ are found by solving $\min_{\theta} \; d \bigl(P, \hat P_\theta\bigr)$.
In practice, this can be solved using any derivative-free optimization algorithm, such as Bayesian optimization or evolutionary algorithms.
Note that one need not be limited to a single prompt; optimization can be solved using multiple templates or sets of agents with different values for the same template.

\paragraph{Measuring Performance.}  
Let $\hat P_{\bm{\theta^{\star}}}$ denote the LLM's distribution of responses under $\bm{\theta^{\star}}$, and $\hat P_0$ the baseline LLM off-the-shelf without any additional prompting.
Training loss is $d(P, \hat P_{\bm{\theta^{\star}}})$ and validation loss is defined analogously on the held-out test setting. 

To assess whether $\bm{\theta^{\star}}$ generalizes, we compare predictive fit against the baseline $\hat P_0$.  
This is an appealing reference measure because many distance metrics between distributions (like the KL divergence) are not easily interpretable.
Furthermore, improvement over the baseline provides direct evidence that designing and optimizing \persona{s} is a worthwhile endeavor in the first place.\footnote{Strong baseline performance would simply reflect the LLM's already high off-the-shelf predictive capacity.}

Relative improvement can be computed as a difference-in-distances: $d\bigl(P, \hat P_{0}\bigr) \;-\; d\bigl(P, \hat P_{\bm{\theta^{\star}}}\bigr)$.
A positive difference indicates that the \persona{s} provide better predictive power than the baseline.
When we optimize over a set of candidate \persona{s}, we seek those that yield this positive difference on both the training and testing data---that would suggest that the set generalizes.

The same framework applies across multiple training and evaluation settings, with optimization performed by averaging distances if needed.  
It also allows direct comparison of optimized \persona{s} to alternative \persona{} sets or benchmark models (e.g., game-theoretic equilibria) using the same measurement tools.  
